\begin{ouverture}[college3]{16}{Résistance et loi d’Ohm}{Partie 3 · L’énergie, ses transferts et ses conversions}{Peut-on prévoir le courant à partir d’une tension ?}
\objectif{Mesurer la caractéristique d’un conducteur ohmique.}
\objectif{Exploiter la relation U = R I avec les unités correctes.}
\objectif{Distinguer une loi de modèle d’une règle valable pour tous les dipôles.}
\prerequis{Tension, intensité, proportionnalité et graphiques.}
\end{ouverture}
\pcPage{Le cours}{Résistance et loi d’Ohm}
\phantomsection\label{ess-16}
\essSection{1}{Caractériser un dipôle}
Pour étudier une résistance, on mesure la tension $U$ à ses bornes et l’intensité $I$ qui la traverse pour plusieurs réglages autorisés de la source. Le voltmètre est en dérivation et l’ampèremètre en série. On ne dépasse pas les caractéristiques du composant, qui peut chauffer. Tracer $U$ en fonction de $I$ permet de rechercher une relation.\par
\pcFigure{figures/c16-montage.pdf}{Le courant traverse A et R ; V compare les deux bornes de R. La source est réglable et limitée selon le matériel scolaire.}
\essSection{2}{Reconnaître un conducteur ohmique}
À température suffisamment constante, la caractéristique d’un conducteur ohmique est une droite passant par l’origine. Le quotient $U/I$ est constant pour les points non nuls. Il définit la résistance $R$, en ohms ($\Omega$). Le nom « résistance » désigne à la fois un dipôle et une grandeur : le contexte doit être clair.\par
\clearpage
\essSection{3}{Appliquer la loi d’Ohm}
La loi s’écrit $U=RI$, avec $U$ en V, $R$ en $\Omega$ et $I$ en A. On en déduit $I=U/R$ ou $R=U/I$ lorsque les données le permettent. Convertir les milliampères en ampères avant de calculer. $1\unite{k\Omega}=1\,000\unite\Omega$. Une résistance plus élevée laisse passer un courant plus faible sous une même tension, dans ce modèle.\par
\pcFigure{figures/c16-ohm.pdf}{La droite la plus pentue sur U en fonction de I correspond à la plus grande résistance.}
\essSection{4}{Connaître les limites}
Tous les dipôles ne sont pas ohmiques. Une lampe à filament chauffe et sa caractéristique peut être courbe ; une diode dépend fortement de la polarité et de la tension. Pour mesurer un composant à l’ohmmètre, on le déconnecte de toute alimentation et on suit la notice. Une valeur expérimentale s’interprète avec tolérance et température, sans lui ajouter des décimales injustifiées.\par
\begin{pcEncadre}{method}{Un exemple expliqué}Une résistance de $150\unite\Omega$ est soumise à $4{,}5\unite V$. Alors $I=U/R=4{,}5/150=0{,}030\unite A=30\unite{mA}$. Inversement, si $I=20\unite{mA}=0{,}020\unite A$, la tension vaut $U=150\times0{,}020=3{,}0\unite V$. On vérifie que le composant supporte le fonctionnement demandé.\end{pcEncadre}
\pcPage{Expérimenter et retenir}{Résistance et loi d’Ohm}
\phantomsection\label{activite-16}
\begin{pcEncadre}{method}{Investigation en classe}\textbf{Question.} Le dipôle fourni est-il compatible avec le modèle ohmique ? Réaliser le montage vérifié par le professeur. Mesurer au moins quatre couples $(I;U)$ dans la plage autorisée, sans échauffement excessif. Construire le graphique $U=f(I)$, puis calculer quelques quotients $U/I$. Comparer la régularité à la précision de mesure. \textbf{À rendre :} schéma, tableau en unités SI, graphique et conclusion limitée à la plage étudiée. L’ohmmètre ne s’utilise jamais sur un circuit alimenté.\end{pcEncadre}
\pcSousTitre{Ce que je dois savoir expliquer}\begin{questions}
\item La caractéristique relie tension et intensité dans des conditions précisées.
\item Un conducteur ohmique vérifie U = R I à température constante.
\item La résistance s’exprime en ohms.
\item Une droite passant par l’origine révèle une proportionnalité.
\item La loi d’Ohm ne décrit pas tous les dipôles.
\end{questions}
\begin{pcEncadre}{attention}{Une erreur à éviter}Utiliser $I=30$ pour $30\unite{mA}$ donnerait un résultat mille fois trop grand avec des ohms : écrire d’abord $I=0{,}030\unite A$.\end{pcEncadre}
\pcRepere{cor-activite-16}{Pistes pour l’investigation}\pcGrille
\clearpage\markright{Je démarre}
\rubrique{Je démarre}
\par\noindent\begin{minipage}{\linewidth}
\begin{exercice}[1]{L’unité de R}
\begin{questions}
\item Quelle est l’unité de résistance ?
\item Quel symbole utilise-t-on ?
\end{questions}
\end{exercice}
\end{minipage}
\par\noindent\begin{minipage}{\linewidth}
\begin{exercice}[1]{La relation}
\begin{questions}
\item Écrire la loi d’Ohm.
\item Donner les unités nécessaires.
\end{questions}
\end{exercice}
\end{minipage}
\par\noindent\begin{minipage}{\linewidth}
\begin{exercice}[1]{Convertir}
\begin{questions}
\item Convertir $40\unite{mA}$ en A.
\item Convertir $2{,}2\unite{k\Omega}$ en ohms.
\end{questions}
\end{exercice}
\end{minipage}
\par\noindent\begin{minipage}{\linewidth}
\begin{exercice}[1]{Une tension directe}
\begin{questions}
\item R vaut $100\unite\Omega$ et I vaut $0{,}02\unite A$. Calculer U.
\item Le produit est-il en volts ?
\end{questions}
\end{exercice}
\end{minipage}
\par\noindent\begin{minipage}{\linewidth}
\begin{exercice}[1]{Une droite}
\pcFigureEx{figures/ex-ohm-modele.pdf}{Repère à compléter qualitativement ; aucune valeur de résistance n’est imposée.}
\begin{questions}
\item Tracer qualitativement U en fonction de I pour un conducteur ohmique idéal.
\item Passe-t-elle par l’origine ?
\end{questions}
\end{exercice}
\end{minipage}
\clearpage\markright{Je m’entraîne}
\rubrique{Je m’entraîne}
\par\noindent\begin{minipage}{\linewidth}
\begin{exercice}[2]{Retrouver le courant}
\begin{questions}
\item $U=6\unite V$ et $R=200\unite\Omega$. Calculer I en A puis mA.
\item Vérifier avec U = R I.
\end{questions}
\end{exercice}
\end{minipage}
\par\noindent\begin{minipage}{\linewidth}
\begin{exercice}[2]{Retrouver R}
\begin{questions}
\item $U=4{,}0\unite V$ pour $I=20\unite{mA}$. Trouver R.
\item Pourquoi la réponse 0,2 ohm serait-elle incorrecte ?
\end{questions}
\end{exercice}
\end{minipage}
\par\noindent\begin{minipage}{\linewidth}
\begin{exercice}[2]{Une table cohérente}
\pcFigureEx{figures/ex-ohm-150.pdf}{}
\begin{questions}
\item On mesure $(I;U)$ : $(0{,}01;1{,}5)$, $(0{,}02;3)$, $(0{,}03;4{,}5)$, en A et V. Trouver R.
\item Prévoir U pour 0,04 A.
\end{questions}
\end{exercice}
\end{minipage}
\par\noindent\begin{minipage}{\linewidth}
\begin{exercice}[2]{Comparer sous la même tension}
\begin{questions}
\item Sous $6\unite V$, calculer I pour $100\unite\Omega$ puis $300\unite\Omega$.
\item Quelle résistance laisse passer le plus grand courant ?
\end{questions}
\end{exercice}
\end{minipage}
\par\noindent\begin{minipage}{\linewidth}
\begin{exercice}[2]{Doubler U}
\begin{questions}
\item Pour R constante, I vaut $15\unite{mA}$ sous $3\unite V$. Prévoir I sous $6\unite V$.
\item Quelle hypothèse est essentielle ?
\end{questions}
\end{exercice}
\end{minipage}
\clearpage\markright{Je raisonne}
\rubrique{Je raisonne}
\par\noindent\begin{minipage}{\linewidth}
\begin{exercice}[3]{Une lampe courbe}
\pcFigureEx{figures/ex-lampe-courbe.pdf}{Exemple de caractéristique non linéaire dans une plage donnée.}
\begin{questions}
\item La caractéristique d’une lampe ne forme pas une droite. Peut-on utiliser une seule valeur de R pour tous ses fonctionnements ?
\item Quelle méthode reste possible ?
\end{questions}
\end{exercice}
\end{minipage}
\par\noindent\begin{minipage}{\linewidth}
\begin{exercice}[3]{Un point suspect}
\pcFigureEx{figures/ex-ohm-point-suspect.pdf}{Droite : modèle de 100 ohms ; croix rouge : relevé à vérifier.}
\begin{questions}
\item Un dipôle de $100\unite\Omega$ donne $2\unite V$ pour $0{,}020\unite A$, mais un relevé indique $5\unite V$ pour $0{,}030\unite A$. Quelle valeur attend-on au second point ?
\item Que vérifier ?
\end{questions}
\end{exercice}
\end{minipage}
\par\noindent\begin{minipage}{\linewidth}
\begin{exercice}[3]{L’ohmmètre}
\begin{questions}
\item Pourquoi déconnecter toute alimentation avant une mesure à l’ohmmètre ?
\item Pourquoi isoler si nécessaire la résistance des autres branches ?
\end{questions}
\end{exercice}
\end{minipage}
\par\noindent\begin{minipage}{\linewidth}
\begin{exercice}[3]{Une tolérance fournie}
\begin{questions}
\item Une résistance annoncée $220\unite\Omega$ admet ici une plage de $209$ à $231\unite\Omega$. Une mesure de $225\unite\Omega$ est-elle compatible ?
\item Sous $4{,}5\unite V$, calculer le courant avec la valeur mesurée.
\end{questions}
\end{exercice}
\end{minipage}
\par\noindent\begin{minipage}{\linewidth}
\begin{exercice}[3]{Un choix raisonné}
\begin{questions}
\item Une source scolaire de $6\unite V$ doit alimenter une résistance avec un courant ne dépassant pas $20\unite{mA}$. Quelle résistance minimale théorique ?
\item Une résistance de 330 ohms convient-elle pour le courant, sous réserve de puissance admissible ?
\end{questions}
\end{exercice}
\end{minipage}
\pcPage{Défi du chapitre}{Résistance et loi d’Ohm}
\phantomsection\label{defi-ex-16}
\pcDefiTitre{16}{Identifier une boîte}
\pcFigureEx{figures/ex-ohm-boite.pdf}{}
\begin{questions}
\item Pour une boîte inconnue, on mesure 2 V à 10 mA, 4 V à 20 mA et 6 V à 30 mA. Estimer R et la nature du modèle.
\item Prévoir le courant à 5 V.
\item Pourquoi ne pas poursuivre automatiquement jusqu’à 100 V ?
\end{questions}
\pcGrille
